\documentclass[aspectratio=169,11pt]{beamer}
\usetheme{default}
\setbeamertemplate{navigation symbols}{}
\setbeamertemplate{footline}{%
  \hfill\insertframenumber/\inserttotalframenumber\hspace*{4pt}\vspace{2pt}%
}
\usepackage{lmodern}
\usepackage[most]{tcolorbox}
\usepackage{listings}
\usepackage{tikz}
\usetikzlibrary{arrows.meta,positioning,shapes.geometric,fit,backgrounds}
\usepackage{booktabs}
\usepackage{hyperref}

\definecolor{lcgreen}{HTML}{2E7D32}
\definecolor{lcblue}{HTML}{1565C0}
\definecolor{lcorange}{HTML}{EF6C00}
\definecolor{lcred}{HTML}{C62828}
\definecolor{codebg}{HTML}{F5F5F5}
\definecolor{docgray}{gray}{0.55}

\newtcolorbox{conceptbox}[1][]{colback=yellow!20,colframe=yellow!60!black,title=#1,fonttitle=\bfseries,top=2pt,bottom=2pt,left=4pt,right=4pt}
\newtcolorbox{warningbox}[1][]{colback=red!15,colframe=red!60!black,title=#1,fonttitle=\bfseries,top=2pt,bottom=2pt,left=4pt,right=4pt}
\newtcolorbox{tipbox}[1][]{colback=green!10,colframe=green!50!black,title=#1,fonttitle=\bfseries,top=2pt,bottom=2pt,left=4pt,right=4pt}
\newtcolorbox{codebox}[1][]{colback=blue!10,colframe=blue!40!black,title=#1,fonttitle=\bfseries,top=2pt,bottom=2pt,left=4pt,right=4pt}

\lstset{
  basicstyle=\ttfamily\scriptsize,
  backgroundcolor=\color{codebg},
  breaklines=true,
  frame=single,
  rulecolor=\color{blue!40!black},
  keywordstyle=\color{lcblue}\bfseries,
  stringstyle=\color{lcgreen},
  commentstyle=\color{docgray}\itshape,
  language=Python,
  showstringspaces=false,
  tabsize=2,
  xleftmargin=2pt,
  xrightmargin=2pt,
  aboveskip=2pt,
  belowskip=2pt,
}

\newcommand{\docref}[1]{{\scriptsize\color{docgray}[Docs: #1]}}

\title{LangGraph v1.0.x}
\subtitle{Section 00: Overview \& Why LangGraph}
\author{Agentic AI Foundations}
\date{March 2026}

\begin{document}

\begin{frame}
\titlepage
\end{frame}

\section{Mental Model}
\begin{frame}[t]{LangGraph: Graph-Based Agent Orchestration}
\setlength{\itemsep}{1pt}
\vspace{0.1cm}
\begin{conceptbox}[Core Idea]
LangGraph is a framework for building \textbf{stateful, multi-step agent applications as directed graphs}, where:
\setlength{\itemsep}{1pt}
\begin{itemize}
  \item \textbf{Nodes} = computational steps (LLM calls, tool execution, data processing)
  \item \textbf{Edges} = transitions between nodes (control flow)
  \item \textbf{State} = shared data that flows through the graph
  \item \textbf{Cycles} = looping paths for agentic decision-making
\end{itemize}
\end{conceptbox}

\vspace{0.15cm}

Think of LangGraph as \textit{``orchestrating conversations as graphs''}-not linear chains, but flexible, cyclic workflows.
\end{frame}

\section{The Problem}
\begin{frame}[t]{Why Chains Aren't Enough}
\setlength{\itemsep}{1pt}
\vspace{0.1cm}
\begin{warningbox}[Limitations of Linear Pipelines]
\setlength{\itemsep}{1pt}
\begin{itemize}
  \item \textbf{No cycles:} can't loop back (agents need to reason -> act → reason again)
  \item \textbf{Stateless:} each step loses context of previous decisions
  \item \textbf{No branching:} can't route to different paths based on conditions
  \item \textbf{No human-in-the-loop:} no pause points for human approval
  \item \textbf{Hard to debug:} linear flow obscures agent decision logic
\end{itemize}
\end{warningbox}

\vspace{0.15cm}

\begin{tipbox}[Example: ReAct Agent]
ReAct (Reasoning + Acting) is inherently cyclic: \textit{Thought} -> \textit{Action} → \textit{Observation} → back to \textit{Thought}. Chains can fake this with structured outputs, but can't truly express the loop.
\end{tipbox}
\end{frame}

\section{What is LangGraph}
\begin{frame}[t]{LangGraph: Definition \& Scope}
\setlength{\itemsep}{1pt}
\vspace{0.1cm}
\textbf{LangGraph} is a low-level library for orchestrating agentic workflows as directed graphs.

\vspace{0.15cm}

\begin{codebox}[What You Build]
\setlength{\itemsep}{1pt}
\begin{itemize}
  \item Agent loops (ReAct, ReWOO, langgraph-agent-patterns)
  \item Multi-turn conversations with state
  \item Workflows with branching, looping, fan-out/fan-in
  \item Persistent, resumable applications
  \item Human approval workflows
\end{itemize}
\end{codebox}

\vspace{0.15cm}

\begin{tipbox}[Philosophy]
LangGraph = \textbf{low-level} orchestration primitives. Build once, use a hundred different ways. Pairs with LangChain (runnable, prompt, tool abstractions).
\end{tipbox}
\end{frame}

\section{LangGraph vs LangChain}
\begin{frame}[t]{LangGraph vs LangChain}
\setlength{\itemsep}{1pt}
\vspace{0.1cm}

\begin{tabular}{p{4cm}p{5.5cm}p{5.5cm}}
\toprule
\textbf{Aspect} & \textbf{LangChain} & \textbf{LangGraph} \\
\midrule
\textbf{Level} & High-level convenience & Low-level orchestration \\
\textbf{Use case} & Quick chains, RAG, prompt templates & Agents, workflows, cycles \\
\textbf{State} & Implicit (context window) & Explicit, typed, persistent \\
\textbf{Cycles} & No native support & Core feature \\
\textbf{Control flow} & Sequential (mostly) & Explicit routing \& branching \\
\textbf{Debugging} & Callbacks, logging & Introspectable graphs \\
\bottomrule
\end{tabular}

\vspace{0.15cm}

\begin{conceptbox}[The Relationship]
LangGraph \textit{uses} LangChain primitives (runnables, tools, chat models). But LangGraph is the orchestration layer for \textbf{multi-step, interactive workflows}.
\end{conceptbox}
\end{frame}

\section{Key Value Props}
\begin{frame}[t]{Key Value Propositions}
\setlength{\itemsep}{1pt}
\vspace{0.1cm}

\begin{tipbox}[5 Pillars of LangGraph]
\setlength{\itemsep}{1pt}
\begin{itemize}
  \item \textbf{Cycles:} Loop back for agentic reasoning (agent patterns)
  \item \textbf{State:} Typed, persistent, multi-step workflows
  \item \textbf{Persistence:} Save/resume execution (checkpointing)
  \item \textbf{Human-in-the-loop:} Pause for human approval or input
  \item \textbf{Streaming:} Stream token-by-token output without loading full response
\end{itemize}
\end{tipbox}

\vspace{0.15cm}

\begin{conceptbox}[Why This Matters]
These features unlock use cases like \textit{multi-step planning}, \textit{tool-use agents}, \textit{approval workflows}, \textit{long-running tasks}, and \textit{complex reasoning chains}.
\end{conceptbox}
\end{frame}

\section{Graph Mental Model}
\begin{frame}[t,shrink=5]{The Graph Mental Model: Nodes, Edges, State}
\setlength{\itemsep}{0.5pt}
\vspace{0.03cm}

{\footnotesize
\begin{tabular}{p{2.6cm}p{5.8cm}}
\toprule
\textbf{Comp.} & \textbf{Definition} \\
\midrule
\textbf{Nodes} & Compute; return updates \\
\textbf{Edges} & Routes between nodes \\
\textbf{State} & Shared data \\
\textbf{Cond. Edge} & Routes on conditions \\
\textbf{START/END} & Entry/exit \\
\bottomrule
\end{tabular}

Nodes->partial updates; graph merges. Analogy: nodes=intersections, edges=roads.
}
\end{frame}

\section{Ecosystem}
\begin{frame}[t,shrink=3]{LangGraph Ecosystem}
\vspace{0.05cm}

{\footnotesize
\begin{conceptbox}[Stack]
\setlength{\itemsep}{0.5pt}
\begin{itemize}
  \item \textbf{Core:} StateGraph, compilation, invoke
  \item \textbf{Checkpoint:} Persistence (save/resume)
  \item \textbf{Platform:} Cloud deployment, jobs
  \item \textbf{LangSmith:} Monitoring, evaluation
\end{itemize}
\end{conceptbox}
}
\end{frame}

\section{Installation}
\begin{frame}[fragile,t]{Installation \& Setup}
\vspace{0.1cm}

\begin{codebox}[pip install]
\begin{lstlisting}
pip install langgraph langgraph-checkpoint
pip install langchain langchain-community
pip install python-dotenv
\end{lstlisting}
\end{codebox}

\vspace{0.1cm}

\begin{tipbox}[API Keys]
Set environment variables for LLM providers:
\begin{lstlisting}
export OPENAI_API_KEY="sk-..."
export ANTHROPIC_API_KEY="sk-ant-..."
\end{lstlisting}
\end{tipbox}

\vspace{0.1cm}

\begin{conceptbox}[Verify Installation]
\begin{lstlisting}
python -c "import langgraph; print(langgraph.__version__)"
# Output: 1.0.x
\end{lstlisting}
\end{conceptbox}
\end{frame}

\section{Your First Program}
\begin{frame}[fragile,t,shrink=3]{Your First LangGraph Program (Minimal)}
\vspace{0.02cm}

\begin{codebox}[Minimal Chatbot]
\begin{lstlisting}
from langgraph.graph import StateGraph, START, END
from typing_extensions import TypedDict

class State(TypedDict):
    messages: list

def chatbot(state):
    return {"messages": state["messages"] + ["Hi"]}

graph = StateGraph(State)
graph.add_node("chatbot", chatbot)
graph.add_edge(START, "chatbot")
graph.add_edge("chatbot", END)
compiled = graph.compile()
result = compiled.invoke({"messages": []})
\end{lstlisting}
\end{codebox}

{\footnotesize Pattern: Define -> Add → Compile → Invoke}
\end{frame}

\section{Package Architecture}
\begin{frame}[t,shrink=5]{Package Architecture}
\setlength{\itemsep}{0.5pt}
\vspace{0.03cm}

\begin{tabular}{p{3.2cm}p{9cm}}
\toprule
\textbf{Package} & \textbf{Purpose} \\
\midrule
\textbf{langgraph} & Core: StateGraph, nodes, edges, compilation \\
\textbf{langgraph.graph} & Graph classes, START/END, builders \\
\textbf{langgraph.checkpoint} & Persistence backends (save/resume) \\
\textbf{langgraph-platform} & Cloud deployment, workers \\
\textbf{langgraph-supervisor} & Multi-agent orchestration \\
\bottomrule
\end{tabular}

\vspace{0.05cm}

\begin{tipbox}[Common Imports]
{\footnotesize \texttt{from langgraph.graph import StateGraph, START, END}}
\end{tipbox}
\end{frame}

\section{Core Concepts}
\begin{frame}[t,shrink=3]{Core Concepts Preview Table}
\vspace{0.05cm}

\begin{tabular}{p{2.5cm}p{2.8cm}p{5.2cm}}
\toprule
\textbf{Concept} & \textbf{Type} & \textbf{Description} \\
\midrule
\texttt{State} & TypedDict & Shared data; merged via reducers \\
\texttt{Reducer} & Function & Merges state updates \\
\texttt{Node} & Function & Computes state updates \\
\texttt{Edge} & Connection & Routes node-to-node \\
\texttt{Conditional Edge} & Router & Routes based on state \\
\texttt{START / END} & Special & Entry/exit points \\
\texttt{StateGraph} & Class & Orchestrates graph \\
\texttt{Compiled Graph} & Runnable & Has \texttt{invoke()}, \texttt{stream()} \\
\bottomrule
\end{tabular}

\vspace{0.1cm}

{\scriptsize Each is fundamental. Sections 01-03 dive deep.}
\end{frame}

\section{LangGraph vs Alternatives}
\begin{frame}[t,shrink=5]{LangGraph vs Alternatives}
\vspace{0.05cm}

{\footnotesize
\begin{tabular}{p{1.6cm}p{2.2cm}p{2.2cm}p{2.2cm}p{2.2cm}}
\toprule
\textbf{} & \textbf{LangGraph} & \textbf{AutoGen} & \textbf{CrewAI} & \textbf{Raw} \\
\midrule
\textbf{State} & Typed & Implicit & Implicit & Ad-hoc \\
\textbf{Cycles} & Native & Yes & Yes & Manual \\
\textbf{Persist} & Yes & No & No & No \\
\textbf{Debug} & Great & Good & Good & Hard \\
\textbf{Flexibility} & Very high & Med & Med & Highest \\
\bottomrule
\end{tabular}
}

\vspace{0.1cm}

\begin{tipbox}[When to Use]
\setlength{\itemsep}{1pt}
{\footnotesize
\begin{itemize}
  \item Fine control over flow, persistence, debugging
\end{itemize}
}
\end{tipbox}
\end{frame}

\section{Learning Path}
\begin{frame}[t,shrink=3]{Learning Path Ahead}
\setlength{\itemsep}{0.5pt}
\vspace{0.05cm}

{\footnotesize
\begin{tipbox}[Sections 00-03]
\setlength{\itemsep}{0.5pt}
\begin{itemize}
  \item \textbf{Sec 00:} Overview
  \item \textbf{Sec 01:} State, nodes, edges
  \item \textbf{Sec 02:} Building graphs
  \item \textbf{Sec 03:} Conditional routing
\end{itemize}
\end{tipbox}

\vspace{0.05cm}

Beyond: persistence, tools, streaming, deployment.
}
\end{frame}

\section{Summary \& Cheat Sheet}
\begin{frame}[t,shrink=4]{Summary \& Cheat Sheet}
\setlength{\itemsep}{0.5pt}
\vspace{0.05cm}

{\footnotesize
\begin{conceptbox}[Key Points]
\setlength{\itemsep}{0.5pt}
\begin{itemize}
  \item Graph-based orchestration (cycles, state, persistence)
  \item Chains can't loop; graphs can
  \item Core: nodes, edges, state
  \item Pattern: define -> add → compile → invoke
\end{itemize}
\end{conceptbox}

Sequential = LangChain. Looping/agentic = LangGraph.
}
\end{frame}

\section{Self-Check}
\begin{frame}[t,shrink=3]{Self-Check Questions}
\setlength{\itemsep}{0pt}
\vspace{0.05cm}

{\footnotesize
\begin{enumerate}
\setlength{\itemsep}{0.5pt}
  \item Why can't linear chains express agent loops?
  \item What are the three core components of LangGraph?
  \item What is a reducer and why does state need one?
  \item How does LangGraph differ from LangChain?
  \item Name two value propositions.
  \item What is the basic workflow to create and run?
  \item When would you choose LangGraph over raw code?
\end{enumerate}

\vspace{0.05cm}

\textbf{Quick Answers:} (1) Chains are linear. (2) Nodes, edges, state. (3) Merges updates; tracks state. (4) Low-level vs high-level. (5) Cycles, persistence. (6) Define -> add → compile → invoke. (7) Complex workflows.
}
\end{frame}

\end{document}
